\chapter{经典模型简介}
\label{chap:classic-models-overview}

同一批房屋成交记录，可以用来预测价格，也可以按预设的价格档位作分类；如果没有预设类别，还可以寻找相似房屋形成的群组，或把许多属性压缩成少数坐标。这些任务要交付的结果不同，却经常使用相近的办法：给特征加权、参考附近样本、按条件划分数据，或用概率分布解释观测。为什么同一种建模思路能够用于不同任务，换了任务以后又需要改变什么？本章为后续章节建立这条联系。

\section{本部分的范围与组织}
\label{sec:classic-overview-scope}

本书用“经典模型”组织回归、分类、聚类、降维及集成学习中的常见方法。这里的“经典”表示已经形成稳定建模主干的方法，不按发表年代划定边界。神经网络的结构与训练由第\ref{part:neural-network-models}部分展开，随机变量之间的图结构及其推断由第\ref{part:probabilistic-models}部分系统讨论。概率建模本身仍贯穿本部分：贝叶斯回归、概率分类和隐变量模型都属于这里需要理解的内容。

第\ref{chap:regression}章根据输入预测数值响应，第\ref{chap:classification}章根据已知类别学习决策；去掉类别答案后，第\ref{chap:clustering}章讨论怎样发现群组，第\ref{chap:dimensionality-reduction}章讨论怎样构造较少的坐标来保留所需信息。集成学习则把多个模型组合起来，可以继续服务于回归或分类等任务。因而，集成描述的是模型的组合方式，与前面四类任务的划分依据不同。

第\ref{chap:machine-learning-models}章已经建立模型、输出和学习准则的基本概念，第\ref{chap:statistical-learning-framework}章说明从训练数据选择模型与评价新样本表现的区别。本章承接这些基础，集中观察经典方法怎样表示规律。具体任务章节再说明参数怎样学习、结果怎样使用，以及什么证据能够支持相应结论。

\section{经典模型的建模思路}
\label{sec:classic-overview-modeling}

以房屋数据为例，预测器可以假定面积等属性共同形成一个全局趋势，也可以主要参考相似房屋，或针对不同条件学习不同的局部规则。概率模型进一步说明相同输入下结果可能怎样变化，集成方法则组合多个成员的判断。这些思路允许交叉：一个模型可以同时使用线性关系和概率分布，也可以把树作为集成成员。下面沿着它们各自回答的问题，说明共享结构怎样随任务改变。

\subsection{基于线性关系的建模}
\label{sec:classic-overview-linear}

预测房价时，一个直接的起点是给面积、楼龄等特征分别赋予权重，再相加得到预测。设输入为$\bm{x}\in\mathbb{R}^{d}$，权重为$\bm{w}\in\mathbb{R}^{d}$，偏置为$b\in\mathbb{R}$，这个计算写成
\begin{equation}
  g(\bm{x})=\bm{w}^{\top}\bm{x}+b.
  \label{eq:classic-overview-linear-score}
\end{equation}
在线性回归中，$g(\bm{x})$可以直接作为价格预测；在二分类中，它可以作为类别得分，再按阈值给出类别。加权求和的结构相同，但回归系数需要贴合数值响应，分类权重需要帮助区分类别。得分也可以经适当转换得到类别概率，这就与后面的概率建模相接。

当目的从预测变成降维时，加权求和仍然有用，只是输出改为少数新的坐标。主成分分析（principal component analysis, PCA）选择若干投影方向，让保留的坐标描述数据的主要变化；这里学习方向的依据是输入自身的变化，而不是房价或类别标签。第\ref{sec:dim-pca}节会说明这一目标与重构误差的关系。同样的线性运算并不意味着相同的学习问题。

如果大户型和小户型的面积增量对应不同的价格变化，或面积的作用取决于楼龄，原始特征的加权和就可能过于受限。可以先加入面积平方、面积与楼龄的乘积等新特征，再学习它们的权重。这种\term{特征映射}（feature mapping）把原来的输入改写成另一组可供计算的特征，使简单的线性组合也能表达原始输入中的非线性关系。线性因而要相对于所用的特征来判断；第\ref{sec:regression-basis}节从基函数展开说明这种扩展。

显式增加特征会扩大表示规模。若方法只通过特征内积使用映射后的输入，就可以直接计算这个内积，形成\term{核方法}（kernel method）。这种内积替换需要合法的核函数，也需要明确核矩阵怎样构造和使用。正定性的证明和Mercer谱展开解释内积表示为何成立，初读时可以先把握定义、条件和常用核的选择，再按需要阅读证明。

图\ref{fig:classification-kernel-map}用一个二维分类例子展示特征映射$\bm{\phi}(\bm{x})=(x_1,x_1^2+x_2^2)^{\top}$。输入空间中的圆$x_1^2+x_2^2=1$，在新坐标中变成直线$\phi_2=1$；同一得分$g(\bm{x})=1-x_1^2-x_2^2$于是可以写成$g(\bm{x})=(0,-1)\bm{\phi}(\bm{x})+1$。改变的是输入的表示，原来位于圆内与圆外的样本不需要改变类别。这幅图给定了一条可行边界，用于解释几何对应，不把它当作训练后的最大间隔解。

\begin{figure}[htbp]
  \centering
  % 教学构造：phi(x)=(x_1,x_1^2+x_2^2)，给定得分 g(x)=1-phi_2(x)。
% 正类：(−0.5,0)、(0,0.4)、(0.4,0.3)；
% 负类：(−1.4,0)、(0,−1.6)、(1.2,0.9)、(0.9,−1.2)。
\begin{tikzpicture}[x=1cm,y=1cm,font=\small]
  \node[font=\heiti\small] at (1.8,4.45) {输入空间：圆形边界};
  \node[font=\heiti\small] at (7.8,4.45) {特征空间：线性边界};
  \begin{scope}[xshift=1.8cm,yshift=1.8cm]
    \fill[BookBlue!6] (-1.65,-1.65) rectangle (1.65,1.65);
    \fill[BookAmber!12] (0,0) circle (1);
    \draw[book arrow] (-1.75,0) -- (1.85,0) node[right] {$x_1$};
    \draw[book arrow] (0,-1.75) -- (0,1.85) node[above] {$x_2$};
    \draw[BookInk,thick] (0,0) circle (1);
    \foreach \x/\y in {-0.5/0,0/0.4,0.4/0.3}
      \fill[BookAmber] (\x-0.065,\y-0.065) rectangle ++(0.13,0.13);
    \foreach \x/\y in {-1.4/0,0/-1.6,1.2/0.9,0.9/-1.2}
      \fill[BookBlue] (\x,\y) circle (2.1pt);
    \node[anchor=north west,font=\footnotesize] at (0.85,-0.7) {$g=0$};
  \end{scope}
  \draw[book arrow,thick] (4.05,1.95) -- node[above] {$\bm{\phi}$} (5.65,1.95);
  \begin{scope}[xshift=7.8cm,yshift=0.15cm]
    \fill[BookAmber!12] (-1.65,0) rectangle (1.65,1);
    \fill[BookBlue!6] (-1.65,1) rectangle (1.65,3.1);
    \draw[book arrow] (-1.75,0) -- (1.85,0) node[right] {$\phi_1$};
    \draw[book arrow] (0,0) -- (0,3.5) node[above] {$\phi_2$};
    \draw[BookInk,thick] (-1.65,1) -- (1.65,1);
    \node[anchor=south east,font=\footnotesize] at (1.65,1) {$g=0$};
    \node[left,font=\footnotesize] at (0,1) {$1$};
    \foreach \x/\y in {-0.5/0,0/0.4,0.4/0.3} {
      \pgfmathsetmacro{\radiusSquared}{\x*\x+\y*\y}
      \fill[BookAmber] (\x-0.065,\radiusSquared-0.065) rectangle ++(0.13,0.13);
    }
    \foreach \x/\y in {-1.4/0,0/-1.6,1.2/0.9,0.9/-1.2} {
      \pgfmathsetmacro{\radiusSquared}{\x*\x+\y*\y}
      \fill[BookBlue] (\x,\radiusSquared) circle (2.1pt);
    }
  \end{scope}
  \fill[BookAmber] (2.1,-0.55) rectangle ++(0.13,0.13);
  \node[right] at (2.35,-0.485) {正类};
  \fill[BookBlue] (4.7,-0.485) circle (2.1pt);
  \node[right] at (4.9,-0.485) {负类};
  \node[font=\footnotesize] at (7.65,-0.485) {$g=1-x_1^2-x_2^2$};
\end{tikzpicture}

  \caption{非线性边界在特征映射前后的对应，两个面板的坐标均无量纲。正类方点为$(-0.5,0)$、$(0,0.4)$、$(0.4,0.3)$，负类圆点为$(-1.4,0)$、$(0,-1.6)$、$(1.2,0.9)$、$(0.9,-1.2)$。右图逐点使用$\phi_1=x_1$、$\phi_2=x_1^2+x_2^2$；水平线$\phi_2=1$对应左图半径为1的圆。数据与边界均为教学构造。}
  \label{fig:classification-kernel-map}
\end{figure}

\subsubsection{特征映射与核技巧}
\label{sec:classic-overview-kernel-trick}
\label{sec:classification-svm-kernel-trick}

设$\bm{\phi}:\mathbb{R}^d\to\mathcal{H}$把输入映射到一个Hilbert空间，即带有完备内积结构、可以讨论长度与正交的特征空间。若能直接计算
\begin{equation}
  K(\bm{x},\bm{z})
  =\langle\bm{\phi}(\bm{x}),\bm{\phi}(\bm{z})\rangle_{\mathcal{H}},
  \label{eq:classic-overview-kernel}
\end{equation}
就可以用这个函数代替显式特征内积，而不先生成$\bm{\phi}(\bm{x})$。这种替换称为\term{核技巧}（kernel trick）。它利用的是模型计算只依赖内积的结构，不要求把无限维向量实际存入内存。

\subsubsection{正定条件与特征空间}
\label{sec:classic-overview-kernel-validity}
\label{sec:classification-svm-mercer}

相似度函数只有满足内积核的条件，才能支持上述替换。有限核矩阵的半正定性刻画一般的内积核；在紧空间、连续性和适当测度条件下，Mercer理论进一步给出正交特征的谱展开。两种结论的条件不同，不能把Mercer定理的拓扑假设当成所有核方法的前提\scite{mercer1909functions}

先允许输入来自任意集合$\mathcal{X}$。以下定理中的$x,z$表示一般输入对象，不要求是数值向量。一个实对称函数$K:\mathcal{X}\times\mathcal{X}\to\mathbb{R}$称为\term{正定核}（positive definite kernel），是指对任意有限输入$x_1,\ldots,x_m$和实系数$a_1,\ldots,a_m$都有
\[
  \sum_{i,j=1}^m a_i a_j K(x_i,x_j)\ge0.
\]
也就是说，对应的\term{Gram矩阵}（Gram matrix）$\bm{K}_{i,j}=K(x_i,x_j)$半正定。这里“正定核”沿用核理论的通常名称，允许Gram矩阵奇异；要求互异输入与非零系数给出严格正值时，才称为严格正定核。

\begin{theorem}[正定核与Hilbert特征映射]
  \label{thm:classification-svm-kernel-characterization}
  对任意集合$\mathcal{X}$上的实对称函数$K$，以下条件等价：
  \begin{enumerate}
    \item 任意有限输入所形成的Gram矩阵均为半正定矩阵。
    \item 存在实Hilbert空间$\mathcal{H}$与映射$\bm{\phi}:\mathcal{X}\to\mathcal{H}$，使$K(x,z)=\langle\bm{\phi}(x),\bm{\phi}(z)\rangle_{\mathcal{H}}$。
  \end{enumerate}
\end{theorem}
\begin{proof}
若特征映射存在，则
\[
  \begin{aligned}
    \bm{a}^{\top}\bm{K}\bm{a}
      &=\sum_{i,j}a_i a_j
         \langle\bm{\phi}(x_i),\bm{\phi}(x_j)\rangle_{\mathcal{H}}\\
      &=\left\lVert\sum_i a_i\bm{\phi}(x_i)\right\rVert_{\mathcal{H}}^2
        \ge0,
  \end{aligned}
\]
所以内积表示必然给出半正定Gram矩阵。

反过来，假设所有有限Gram矩阵都半正定。为每个$x\in\mathcal{X}$引入一个形式基向量$\bm{e}_x$，构造所有有限线性组合组成的向量空间$\mathcal{V}_0$。若$\bm{u}=\sum_i a_i\bm{e}_{x_i}$、$\bm{v}=\sum_j b_j\bm{e}_{z_j}$，定义双线性形式
\[
  \langle\bm{u},\bm{v}\rangle_0
  =\sum_{i,j}a_i b_j K(x_i,z_j).
\]
形式基向量按输入索引，重复输入的系数先合并，所以定义不依赖有限和的书写方式。对称性来自$K$，非负性来自有限Gram矩阵假设。唯一尚未满足的内积条件是：一个非零形式向量也可能具有零长度。

为去掉这些零长度方向，先证明它们与一切向量正交。对任意$t\in\mathbb{R}$，有
\[
  0\le\langle\bm{u}+t\bm{v},\bm{u}+t\bm{v}\rangle_0.
\]
把右侧看成关于$t$的二次多项式可得Cauchy--Schwarz不等式$|\langle\bm{u},\bm{v}\rangle_0|^2\le\langle\bm{u},\bm{u}\rangle_0\langle\bm{v},\bm{v}\rangle_0$；若其中一个平方长度为零，交叉内积必须为零。于是零空间
\[
  \mathcal{N}
  =\{\bm{u}\in\mathcal{V}_0:
       \langle\bm{u},\bm{u}\rangle_0=0\}
\]
是线性子空间，且在商空间$\mathcal{V}_0/\mathcal{N}$上，$\langle[\bm{u}],[\bm{v}]\rangle=\langle\bm{u},\bm{v}\rangle_0$与代表元的选择无关。商掉零方向以后，该形式成为真正的内积。对这个内积空间作完备化，得到Hilbert空间$\mathcal{H}$。

定义$\bm{\phi}(x)=[\bm{e}_x]$，并把这些等价类视为完备化空间中的元素，便有
\[
  \langle\bm{\phi}(x),\bm{\phi}(z)\rangle_{\mathcal{H}}
  =\langle\bm{e}_x,\bm{e}_z\rangle_0=K(x,z).
\]
这就直接从有限Gram半正定性构造出了所需特征空间。
\end{proof}

该构造还可以把$\mathcal{H}$中的元素解释为函数。对$\bm{u}\in\mathcal{H}$，令$f_{\bm{u}}(x)=\langle\bm{u},\bm{\phi}(x)\rangle_{\mathcal{H}}$。若这个函数处处为零，$\bm{u}$就与特征向量的稠密线性张成空间正交，只能是零向量，所以这种函数表示没有丢失信息。记$K_x=K(\,\cdot\,,x)$，对应的函数空间满足
\[
  f(x)=\langle f,K_x\rangle_{\mathcal{H}},
  \qquad
  |f(x)|\le\lVert f\rVert_{\mathcal{H}}\sqrt{K(x,x)}.
\]
从一个函数与核截面的内积中即可恢复该函数的点值，这就是\term{再生核Hilbert空间}（reproducing kernel Hilbert space，RKHS）的“再生”性质。上述一般构造属于Moore--Aronszajn理论；它不依赖紧性、连续性或积分测度\scite{aronszajn1950kernels}

如果希望把这种抽象特征具体写成一列正交坐标，还需要分析一个积分算子的谱。下面给出Mercer谱展开的一个常用现代形式；有限Gram半正定性依旧是核心假设，而空间和测度条件保证展开能够在所有输入点上成立。

\begin{theorem}[连续核的Mercer谱展开]
  \label{thm:classification-svm-mercer}
  设$\mathcal{X}$为紧度量空间，$\mu$为有限Borel测度且具有全支持，即每个非空开集的测度都为正。若$K$为$\mathcal{X}\times\mathcal{X}$上的实对称连续正定核，则积分算子
  \[
    (T_K f)(x)=\int_{\mathcal{X}}K(x,z)f(z)\,\mathrm{d}\mu(z)
  \]
  在$L^2(\mu)$上为紧、自伴、半正定算子。它的正本征值可列为有限或可数序列$\mu_j>0$，相应本征函数$\psi_j$可选为连续且在$L^2(\mu)$中正交归一，并有
  \begin{equation}
    K(x,z)=\sum_j\mu_j\psi_j(x)\psi_j(z).
    \label{eq:classification-svm-mercer-expansion}
  \end{equation}
  展开在$\mathcal{X}\times\mathcal{X}$上绝对且一致收敛；零核对应空和。
\end{theorem}
\begin{proof}
利用上一构造，先把$K$对应的RKHS记为$\mathcal{H}$。核截面满足
\[
  \lVert K_x-K_z\rVert_{\mathcal{H}}^2
  =K(x,x)+K(z,z)-2K(x,z).
\]
由$K$连续可知$x\mapsto K_x$在$\mathcal{H}$中连续，因此每个$f\in\mathcal{H}$也连续。$\mathcal{H}$单位球上的函数一致有界，且由
\[
  |f(x)-f(z)|
  \le\lVert f\rVert_{\mathcal{H}}
     \lVert K_x-K_z\rVert_{\mathcal{H}}
\]
得到等度连续性。Arzel\`a--Ascoli定理保证这个单位球在连续函数的一致范数下相对紧；由于$\mu$有限，嵌入算子$J:\mathcal{H}\to L^2(\mu)$也是紧算子。

全支持条件使$J$单射：连续函数若在某点非零，就在该点的一个开邻域中与零保持正距离，而这个邻域测度为正，所以它不可能在$L^2(\mu)$中等于零。令$S=J^{\ast}J$，其中$J^{\ast}$为伴随算子，则$S$在$\mathcal{H}$上紧、自伴、半正定且核空间为零。紧度量空间具有可数稠密子集，相应核截面张成的空间在$\mathcal{H}$中稠密，因此$\mathcal{H}$可分。由紧自伴算子的谱定理，可以取$\mathcal{H}$的一组正交归一基$\{e_j\}$满足
\[
  Se_j=\mu_j e_j,\qquad \mu_j>0.
\]
若空间有限维，下述求和就是有限和。

伴随的定义与再生性质给出
\[
  (J^{\ast}f)(x)
  =\int_{\mathcal{X}}K(x,z)f(z)\,\mathrm{d}\mu(z),
\]
所以$T_K=JJ^{\ast}$，从而具有上述紧性、自伴性与半正定性。定义$\psi_j=e_j/\sqrt{\mu_j}$，则
\[
  \langle\psi_j,\psi_k\rangle_{L^2(\mu)}
  =\frac{\langle e_j,Se_k\rangle_{\mathcal{H}}}
         {\sqrt{\mu_j\mu_k}}
  =\mathbb{1}\{j=k\},
  \qquad T_K\psi_j=\mu_j\psi_j.
\]
这里$e_j$本来就是连续函数，故$\psi_j$也有连续版本；$JJ^{\ast}$的非零谱与$J^{\ast}J$一致。

对核截面$K_x,K_z$在基$\{e_j\}$中使用Parseval恒等式，并利用$\langle K_x,e_j\rangle_{\mathcal{H}}=e_j(x)$，得到
\[
  K(x,z)
  =\langle K_x,K_z\rangle_{\mathcal{H}}
  =\sum_j e_j(x)e_j(z)
  =\sum_j\mu_j\psi_j(x)\psi_j(z).
\]
特别地，对角尾项
\[
  r_m(x)=K(x,x)-\sum_{j=1}^m e_j(x)^2
\]
是非负连续函数，随$m$单调下降并逐点趋于零。Dini定理将其加强为紧空间上的一致收敛。再由Cauchy--Schwarz不等式，
\[
  \sum_{j>m}|e_j(x)e_j(z)|
  \le\sqrt{r_m(x)r_m(z)}
  \le\sup_{u\in\mathcal{X}}r_m(u)\longrightarrow0.
\]
因此非对角展开也绝对且一致收敛。
\end{proof}

这一谱展开给出了具体的序列特征
\[
  \bm{\phi}(x)
  =\bigl(\sqrt{\mu_1}\psi_1(x),
          \sqrt{\mu_2}\psi_2(x),\ldots\bigr)^{\top}\in\ell^2.
\]
它确实属于平方可和序列空间，因为其平方范数为$K(x,x)<\infty$。Mercer的原始工作研究正负型函数与积分方程的联系；上面采用的是适合核学习表述的紧度量空间版本。

全支持测度在证明中的作用也已清楚：它让连续函数的逐点信息不被$L^2$等价类丢掉。对非紧输入域或不连续核，一般Hilbert特征映射仍可能存在，只是不能直接使用这一版本的一致谱展开。

实际设计核时，不必每次重做谱分解。非负系数加权的核之和，可用特征直和构造；两个核的乘积，可用张量积特征构造，也可由Gram矩阵的Schur乘积定理验证半正定性。若一列正定核逐点收敛到有限值，则任意有限二次型的极限仍非负，极限也是正定核。这些闭包性质提供了可复用的构造规则。检查某一份训练集的Gram矩阵只能认证该有限矩阵；要使核在新样本上也合法，还需要针对整个定义域的证明或已知构造。

\subsubsection{常用核与表示选择}
\label{sec:classic-overview-common-kernels}
\label{sec:classification-svm-common-kernels}

回到$\bm{x},\bm{z}\in\mathbb{R}^d$。常用核之间的区别主要在于怎样度量接近程度、允许哪些特征交互，以及这些选择带来怎样的函数范数。表\ref{tab:classification-svm-kernels}列出五种常见形式，核尺度统一记为$\beta$，多项式次数记为$p$，避免与输入维数$d$混用。

\begin{table}[htbp]
  \centering
  \small
  \begin{tabularx}{\linewidth}{@{}>{\RaggedRight\arraybackslash}p{18mm}>{\RaggedRight\arraybackslash}X>{\RaggedRight\arraybackslash}p{27mm}@{}}
    \toprule
    名称 & $K(\bm{x},\bm{z})$ & 参数与正定性\\
    \midrule
    线性核 & $\bm{x}^{\top}\bm{z}$ & 半正定\\
    多项式核 & $(\beta\bm{x}^{\top}\bm{z}+r)^p$ & $\beta,r\ge0$，正整数$p$\\
    高斯RBF核 & $\exp(-\beta\lVert\bm{x}-\bm{z}\rVert_2^2)$ & $\beta>0$，半正定\\
    Laplacian核 & $\exp(-\beta\lVert\bm{x}-\bm{z}\rVert_1)$ & $\beta>0$，半正定\\
    Sigmoid形式 & $\tanh(\beta\bm{x}^{\top}\bm{z}+r)$ & 一般不保证半正定\\
    \bottomrule
  \end{tabularx}
  \caption{常用核函数与合法性条件。这里的“半正定”要求任意有限输入的Gram矩阵半正定；Sigmoid一行列出的是常见相似性形式，使用前仍需针对定义域与参数核验。}
  \label{tab:classification-svm-kernels}
\end{table}

\term{线性核}（linear kernel）对应$\bm{\phi}(\bm{x})=\bm{x}$，保留原始特征内积。对于线性预测，训练样本的核展开可以聚合为一个权重向量，使用时不必逐一访问这些样本。是否需要额外的非线性交互，应通过任务与验证数据判断；输入维数高本身既不保证线性结构足够，也不意味着非线性核必然更好。

\term{多项式核}（polynomial kernel）把有限阶交互隐式加入模型。以$d=2$、$\beta=1$、$r=0$、$p=2$为例，
\[
  \begin{aligned}
    (\bm{x}^{\top}\bm{z})^2
      &=x_1^2z_1^2+2x_1x_2z_1z_2+x_2^2z_2^2,\\
    \bm{\phi}(\bm{x})
      &=(x_1^2,\sqrt2\,x_1x_2,x_2^2)^{\top}.
  \end{aligned}
\]
交叉项前的$\sqrt2$使特征内积准确等于核值。一般地，令多重指标$\bm{a}=(a_1,\ldots,a_d)$满足$|\bm{a}|=\sum_j a_j=p$，并记$\bm{x}^{\bm{a}}=\prod_j x_j^{a_j}$、$\bm{a}!=\prod_j a_j!$。多项式展开为
\[
  (\bm{x}^{\top}\bm{z})^p
  =\sum_{|\bm{a}|=p}\frac{p!}{\bm{a}!}
     \bm{x}^{\bm{a}}\bm{z}^{\bm{a}}.
\]
所以齐次$p$次核可以使用$\binom{d+p-1}{p}$个单项式坐标；当$\beta>0,r>0$时，非齐次核包括$0$至$p$次项，一种完整单项式表示的维数是$\binom{d+p}{p}$。这些是特征数量，不是样本数量；计算一个核值通常只需一次$d$维内积和幂运算，节省的是显式特征展开成本，完整优化仍要另计。

\term{高斯径向基核}（Gaussian radial basis function kernel，Gaussian RBF kernel）根据欧氏距离赋予局部相似性。其无限维特征可以由指数级数直接看出：
\begin{equation}
  \begin{aligned}
    K(\bm{x},\bm{z})
      ={}&e^{-\beta\lVert\bm{x}\rVert_2^2}
          e^{-\beta\lVert\bm{z}\rVert_2^2}\\
       &\cdot\sum_{p=0}^{\infty}
          \frac{(2\beta)^p}{p!}
          (\bm{x}^{\top}\bm{z})^p.
  \end{aligned}
  \label{eq:classification-svm-rbf-expansion}
\end{equation}
每一阶多项式核都是半正定的，系数非负，两个单点因子可以吸收到特征中；级数逐点收敛，因而也给出了高斯核合法性的证明。把每一阶的单项式展开合并，坐标可写成
\[
  \phi_{\bm{a}}(\bm{x})
  =e^{-\beta\lVert\bm{x}\rVert_2^2}
     \sqrt{\frac{(2\beta)^{|\bm{a}|}}{\bm{a}!}}\,
     \bm{x}^{\bm{a}},
  \qquad \bm{a}\in\mathbb{N}_0^d.
\]
这里$\mathbb{N}_0$为非负整数集合，所有坐标平方之和等于$K(\bm{x},\bm{x})=1$。在整个$\mathbb{R}^d$上，这个核具有无限秩；限制到一批有限训练点以后，实际使用的Gram矩阵仍是$n\times n$。

参数$\beta$控制相似性随距离衰减的速度，也常写成$\beta=1/(2\sigma^2)$。较大$\beta$使不同训练点之间的相似性迅速减小，模型能够形成更局部的变化；较小$\beta$使作用范围更宽。对互异有限输入，高斯Gram矩阵严格正定，因此在不约束范数时，线性方程$\bm{K}\bm{a}=\bm{y}$允许核展开插值给定标签。然而所需系数与函数范数可能很大；固定范数、固定正则化强度或有重复输入且标签冲突时，都不能由“无限维”推出任意插值，更不能由插值推出泛化良好。

\term{Laplacian核}（Laplacian kernel）使用$L^1$距离，且可以分解成各坐标上的乘积：
\[
  \exp(-\beta\lVert\bm{x}-\bm{z}\rVert_1)
  =\prod_{j=1}^d\exp(-\beta|x_j-z_j|).
\]
一维形式的半正定性可由积分特征表示验证：
\[
  e^{-\beta|t|}
  =\int_{\mathbb{R}}e^{\mathrm{i}\omega t}
      \frac{\beta}{\pi(\beta^2+\omega^2)}\,\mathrm{d}\omega.
\]
其中$\mathrm{i}$为虚数单位。对任意有限个输入$x_j$和实系数$a_j$，将有限求和移入积分可得
\[
  \begin{aligned}
    &\sum_{j,k}a_j a_k e^{-\beta|x_j-x_k|}\\
    &\quad=\int_{\mathbb{R}}
      \left|\sum_j a_j e^{\mathrm{i}\omega x_j}\right|^2
      \frac{\beta}{\pi(\beta^2+\omega^2)}\,\mathrm{d}\omega
      \ge0.
  \end{aligned}
\]
多维形式再由核的乘积闭包得到。Laplacian核在零距离附近的变化方式与高斯核不同，可以作为直方图或稀疏特征上的另一种距离选择，但哪一种更适合仍取决于尺度、表示与验证结果。

常称为\term{Sigmoid核}（sigmoid kernel）的形式将内积送入双曲正切函数，形态上接近一个具有非线性激活的神经元，却不自动具有内积核的性质。甚至$\beta=1,r=0$也不能在整个实数域上保证半正定：对$x_1=1,x_2=2$，二阶Gram行列式为
\[
  \tanh(1)\tanh(4)-\tanh(2)^2
  \approx-0.1683<0.
\]
因此“参数取正”并不足以保证合法。某个有限数据集上矩阵没有负特征值，只说明当前矩阵通过检查；若要保留依赖内积核的几何与优化解释，应选用在目标定义域上已确认半正定的核。

对字符串、图与树，核技巧同样适用。\term{字符串核}（string kernel）可以计数共同的子串或子序列，\term{图核}（graph kernel）可以比较局部子结构，\term{树核}（tree kernel）可以比较子树模式。若将每种模式的出现次数作为一个坐标，两对象的模式计数内积就给出了直接的正定构造；动态规划或结构化计数有时能避免把全部坐标显式列出。这样的\term{结构化核}（structured kernel）把输入的组合结构用于相似性计算，适用于序列分析、分子图或句法树等场景，并通过可验证的特征内积保留核方法的几何与优化性质。

\subsubsection{核矩阵的使用与中心化}
\label{sec:classic-overview-kernel-matrix}

在$n$个训练输入上，核矩阵的元素为$\bm{K}_{i,j}=K(\bm{x}_i,\bm{x}_j)$。完整存储需要$O(n^2)$空间和$O(n^2)$次核计算；显式特征可以很大甚至无限维，并不意味着样本之间的计算也消失了。缓存、稀疏化或低秩近似会改变实现成本，近似误差仍须与具体任务一起评价。

有些方法要求在特征空间中先减去训练均值。令$\bm{1}\in\mathbb{R}^{n}$为全一向量、$\bm{I}_n$为单位矩阵，中心化矩阵及中心化核矩阵为
\begin{equation}
  \bm{H}=\bm{I}_n-\frac1n\bm{1}\bm{1}^{\top},
  \qquad \bm{K}_{\mathrm{c}}=\bm{H}\bm{K}\bm{H}.
  \label{eq:classic-overview-kernel-centering}
\end{equation}
左乘$\bm{H}$减去各列均值，右乘$\bm{H}$减去各行均值；展开双边乘积时，还要加回原核矩阵的整体均值。结果等于中心化特征的内积矩阵。只对原始输入作中心化，一般不能代替非线性特征中心化。新样本也应使用训练特征的均值，不能重新用测试数据定义中心。中心化是否必要由具体目标决定，例如核PCA需要它，高斯过程的协方差建模则不因此一律先作这种变换。


核支持向量分类器在特征空间中学习分界，核支持向量回归拟合数值响应，核PCA保留特征空间方差，核K-means则优化特征空间的簇内距离。四者共享核定义的几何，但学习目标不同；相应入口见第\ref{sec:classification-svm-kernels}、第\ref{sec:regression-kernel-svr}、第\ref{sec:dim-kpca}和第\ref{sec:clustering-partition}节。非线性表达来自映射，映射是否适合任务仍需结合数据判断。

\subsection{基于邻域的建模}
\label{sec:classic-overview-neighborhood}

另一种做法是直接寻找与待预测房屋相似的成交记录。\term{邻域}（neighborhood）指按选定距离或相似关系找到的一组附近样本；模型主要利用这些样本的信息，而不必为所有房屋规定同一个全局趋势。$k$近邻（$k$-nearest neighbors, $k$-NN）用最近的$k$个样本构成邻域：回归时可以平均它们的价格，分类时可以让它们的类别投票，也可以给更近的样本更大权重。改变的是聚合的对象和产生输出的规则，寻找邻居的计算则可以复用。

没有响应或类别时，邻域还可以提供样本之间的结构证据。密度聚类考察一个位置附近是否有足够多的样本，再按相应规则连接稠密区域；基于图的方法把样本作为节点，把选定的相似关系作为边。图结构可以用于寻找群组，也可以在少量样本有标签时约束标签传播。用于降维时，则可以要求原本相近的样本在低维表示中仍保持邻近。距离、邻域和图因此连接了预测、聚类与表示学习，但“同一簇”“同一类别”和“相近坐标”仍是不同的要求。

也可以用少量\term{原型}（prototype）代表样本群体，再按到原型的距离作判断。最近质心分类器用类别中心比较输入，K-means则在没有类别标签时同时学习中心和簇分配。两者都比较样本与代表，但前者的类别已知，后者需要寻找分组；原型也不等同于查询点的局部邻居。第\ref{sec:classification-instances-comparison}节与第\ref{sec:clustering-partition}节分别讨论这些使用方式。

这些做法都依赖“相近”的含义。\term{相异度}（dissimilarity）用较小值表示更相近，不一定满足距离度量的全部条件；需要三角不等式的计算不能直接用于任意相异度。房屋在面积上接近，却可能位于不同地区；把面积、楼龄和位置直接混在一个距离中，也可能让数值尺度较大的属性主导邻居选择。表示和度量决定邻域，邻域再影响预测、连通与布局。第\ref{sec:regression-instance}节和第\ref{sec:classification-instances}节比较局部预测，第\ref{sec:clustering-density}节与第\ref{sec:clustering-graph}节讨论群组结构，第\ref{sec:dim-neighbors}节讨论低维布局应保留哪些邻域关系。

\subsection{基于树与规则的建模}
\label{sec:classic-overview-tree-rules}

如果不同房屋适合采用不同的预测规则，可以通过具体条件划分输入。例如，对某个地区的房屋按面积划分，再对其中一组按楼龄划分。\term{决策树}（decision tree）把这些条件组织成逐层分支的结构：输入沿分支到达一个叶节点，再使用该叶的预测。每条路径规定一个区域，不同区域可以采用不同的局部判断。

回归树和分类树共享这套划分结构。采用平方损失时，常数叶的回归树在叶内给出响应均值；分类树可以给出类别或类别比例。因为要拟合的目标不同，选择分裂的准则也随之改变。回归还可以在叶内拟合线性模型，使一片区域既有条件边界，也有局部趋势。这里同时用到了分区和线性关系，说明建模思路可以组合。第\ref{sec:regression-trees}节与第\ref{sec:classification-trees}节分别展开这些训练与预测过程。

树路径也可以写成“如果若干条件成立，就给出某个结论”。\term{规则模型}（rule-based model）直接学习和组织这类条件，规则之间可以重叠，也可能有优先顺序，不必形成一棵完整的树。第\ref{sec:classification-rules}节会说明怎样处理覆盖、冲突和默认类别。浅树或短规则便于检查具体条件，但随着分支或规则数量增加，理解整个决策过程也会变得困难；只凭表示形式不能认定一个模型已经容易解释。

树还可以组织嵌套群组。层次聚类逐步合并或分裂样本组，保留不同粒度下的包含关系；其树描述分组层次，决策树的分支则描述输入条件与预测路径。二者都使用树形结构，但学习准则与输出不同，不能直接共用分类树的分裂规则，详见第\ref{sec:clustering-hierarchy}节。

\subsection{基于概率的建模}
\label{sec:classic-overview-probabilistic}

同样面积、楼龄和位置的房屋，成交价格仍会变化。只输出一个数值，会省去这种变化的描述。\term{概率建模}（probabilistic modeling）为相关变量建立概率分布，用分布说明哪些结果更可能出现，以及结果可能怎样分散。回归可以描述给定输入后的响应分布，也可以对未知系数或函数设置先验，再根据观测形成后验。响应本身的变化与有限数据造成的估计不确定性由此有了不同的建模位置；第\ref{sec:regression-probabilistic}节进一步说明这一差别。

分类关注的则是给定输入后的类别概率。可以直接建立这一条件分布，也可以先描述各类别内部的输入分布及类别出现的比例，再通过贝叶斯公式计算类别概率。前一种思路通向Logistic回归等条件模型，后一种思路通向朴素贝叶斯等生成式分类器。朴素贝叶斯用“给定类别后各特征条件独立”的假设简化联合分布，这个假设不等于说所有房屋属性在总体中互不相关。第\ref{sec:classification-probability}节比较这些分布设定怎样改变训练与决策。

若类别没有给出，可以把群组归属作为\term{隐变量}（latent variable），即参与解释观测、却没有直接记录在数据中的量。高斯混合模型（Gaussian mixture model, GMM）用多个成分分布共同描述输入，并估计一个样本属于各成分的概率；聚类结果的含义依赖这些成分是否适合所要寻找的群组。降维中的隐变量承担另一种职责：它是少量潜在坐标，模型用这些坐标及噪声生成高维观测。概率主成分分析把这一生成关系设为线性，高斯过程隐变量模型则允许非线性映射。第\ref{sec:clustering-gmm}节与第\ref{sec:dim-probabilistic}节会分别说明离散群组和低维坐标的作用。

概率因此可以与前面的表示方式结合。贝叶斯线性回归仍使用线性函数，高斯过程通过核描述函数值之间的协方差，隐变量模型可以把线性或非线性映射写进生成过程。但模型给出了概率，不意味着这些概率已经符合真实数据；分布假设、参数估计与校准仍影响它们的解释。有关概率输出的经验检验可回看第\ref{sec:reliability-output-evaluation}节。

\subsection{基于集成的建模}
\label{sec:classic-overview-ensemble}

同一个预测任务可以训练出不同模型。一棵树可能对训练样本的少量变化很敏感，一个简单线性模型也可能留下系统性的预测误差。\term{集成学习}（ensemble learning）通过训练多个成员并组合它们的输出，改变整体模型的预测方式。成员通常称为\term{基学习器}（base learner）；组合能否带来收益，既取决于成员各自的表现，也取决于它们是否反复犯相同的错误。

针对训练扰动，可以从原训练集有放回地抽取多份数据，通常每份与原训练集的样本数相同，分别训练成员，再平均回归预测或汇总分类投票。这种自助采样聚合称为\term{Bagging}（bootstrap aggregating）：用重采样产生不同成员，用聚合减弱不稳定预测的波动。随机森林还在节点分裂时随机选择候选特征，让各棵树在不同的信息条件下生长。成员共享相近的任务和模型形式，但不需要得到相同的划分；误差相关性仍会限制聚合的收益\scite{breiman1996bagging,breiman2001randomforests}

针对尚未拟合好的部分，可以让新成员接着已有预测训练。\term{Boosting}（提升）逐轮加入成员，使后续学习回应当前整体模型的不足。AdaBoost依据当前分类错误重新加权样本，让分错的样本在下一轮更受重视；梯度提升则根据所选损失确定每一轮应拟合的修正方向，平方损失回归中的修正可以直接表现为当前残差。任务与损失改变后，修正的含义也会改变，所以“每轮拟合残差”不能概括所有Boosting方法。弱学习与提升的理论联系见第\ref{sec:boosting}节，梯度提升的具体构造来自函数空间中的逐步优化\scite{friedman2001gradientboost}

如果线性模型、近邻和树各有可用信息，还可以学习怎样利用它们的预测。\term{Stacking}（堆叠）把基学习器的预测作为另一层模型的输入，由这个\term{元学习器}（meta-learner）学习组合规则。训练组合层时，需要让每个样本对应的基预测来自未使用该样本拟合的模型，例如交叉验证中的折外预测；否则，组合层容易学到基模型在训练样本上的过度乐观表现。Stacking可以混合不同类型的成员，但也增加了训练步骤和预测时需要运行的模型数量\scite{wolpert1992stacked}

\Needspace{18\baselineskip}
\section{本章小结}
\label{sec:classic-overview-summary}

回到房屋数据，不同任务决定要预测价格、判断类别、发现群组还是保留表示；建模思路则决定采用什么关系完成这些目标。一个加权和可以是价格预测、类别得分或低维坐标，一个邻域可以提供响应、标签或结构证据。共享计算需要与任务目标一起理解，不能仅由方法名称推断输出含义。

这些思路也不是相互排斥的选择。线性关系可以借助特征映射与核扩展，局部区域可以包含线性模型，概率分布可以附着在线性或核表示上，多个成员还可以进一步集成。后续章节将以具体任务为依据，说明这些组合怎样学习、怎样评价，以及在什么条件下适用。从第\ref{chap:regression}章开始，数值响应提供了观察这些差异的具体起点。
